% Standalone problem document, generated from the adjacent README.md.
% Compile with XeLaTeX. Mathematical content is maintained in Markdown.
\documentclass[11pt,a4paper]{article}
\usepackage[fontsize=10.5pt]{scrextend}
\usepackage[margin=22mm,headheight=15pt,headsep=9mm,footskip=11mm]{geometry}
\usepackage{fontspec}
\setmainfont{texgyrepagella}[Extension=.otf,UprightFont=*-regular,BoldFont=*-bold,ItalicFont=*-italic,BoldItalicFont=*-bolditalic]
\setsansfont{texgyreheros}[Extension=.otf,UprightFont=*-regular,BoldFont=*-bold,ItalicFont=*-italic,BoldItalicFont=*-bolditalic]
\setmonofont{texgyrecursor}[Extension=.otf,UprightFont=*-regular,BoldFont=*-bold,ItalicFont=*-italic,BoldItalicFont=*-bolditalic,Scale=MatchLowercase]
\usepackage{amsmath,amssymb}
\usepackage{unicode-math}
\setmathfont{texgyrepagella-math.otf}
\usepackage{xcolor,microtype,enumitem,needspace,fancyhdr}
\usepackage{bookmark}
\definecolor{ink}{HTML}{17344B}
\definecolor{accent}{HTML}{197D80}
\definecolor{muted}{HTML}{596773}
\hypersetup{colorlinks=true,linkcolor=accent,urlcolor=accent,pdftitle={RE-03 - Optimal
matvec query complexity of HODLR
approximation},pdfauthor={Open Problems in Numerical Linear Algebra}}
\urlstyle{same}
\setlength{\parindent}{0pt}
\setlength{\parskip}{5pt plus 1pt minus 1pt}
\linespread{1.04}
\setlength{\emergencystretch}{3em}
\widowpenalty=10000
\clubpenalty=10000
\displaywidowpenalty=10000
\allowdisplaybreaks[1]
\setlist{nosep,leftmargin=1.5em}
\providecommand{\tightlist}{\setlength{\itemsep}{0pt}\setlength{\parskip}{0pt}}
\providecommand{\pandocbounded}[1]{#1}
\pagestyle{fancy}
\fancyhf{}
\fancyhead[L]{\sffamily\footnotesize\color{muted}OPEN PROBLEMS IN NUMERICAL LINEAR ALGEBRA}
\fancyhead[R]{\sffamily\bfseries\footnotesize\color{accent}RE-03}
\fancyfoot[L]{\parbox[b]{0.9\textwidth}{\sffamily\fontsize{8}{10}\selectfont\color{muted}Editorial ratings; a bounded search cannot certify that a problem remains unsolved.\\Literature check: 2026-09-13}}
\fancyfoot[R]{\sffamily\footnotesize\color{muted}\thepage}
\renewcommand{\headrulewidth}{0.3pt}
\renewcommand{\headrule}{\hbox to\headwidth{\color{accent}\leaders\hrule height \headrulewidth\hfill}}
\makeatletter
\renewcommand{\section}{\Needspace{5\baselineskip}\@startsection{section}{1}{0pt}{12pt plus 2pt minus 2pt}{4pt}{\sffamily\bfseries\large\color{ink}}}
\renewcommand{\subsection}{\Needspace{4\baselineskip}\@startsection{subsection}{2}{0pt}{9pt plus 1pt minus 1pt}{3pt}{\sffamily\bfseries\normalsize\color{ink}}}
\makeatother
\setcounter{secnumdepth}{0}
\begin{document}
{\sffamily\small\color{muted}Randomized and low-rank approximation\par}
\vspace{3pt}
{\raggedright\hyphenpenalty=10000\exhyphenpenalty=10000\sffamily\bfseries\fontsize{21}{24}\selectfont\color{ink}Optimal
matvec query complexity of HODLR approximation\par}
\vspace{7pt}
{\sffamily\small\textbf{Difficulty:} challenging\quad\textbf{Importance:} interesting
to the community\par}
{\sffamily\small\bfseries\color{accent}Status: Partially resolved\par}
\vspace{4pt}
{\color{accent}\hrule height 0.8pt}
\vspace{6pt}
\textbf{Rating rationale:} Challenging because matching adaptive-query
lower and upper bounds must cover depth and accuracy jointly; community
impact is the cost of hierarchical matrix compression.

\textbf{Area:} information complexity of hierarchical matrix
approximation

\subsection{Partial resolution --- 13 September
2026}\label{partial-resolution-13-september-2026}

Sidney Holden (Center for Computational Biology, Flatiron Institute,
Simons Foundation) proves the following bounds in
\href{https://github.com/ajt60gaibb/OpenProblemsInNLA/blob/main/references/holden-re03-2026-09-13/submission/manuscript/re03_extended_results.pdf}{Theorem
1.1 of the continuation manuscript}, for all original parameters and
universal constants \(c,C>0\):

\[c\min\{n,kL/\varepsilon+k/\varepsilon^2\}
\le q_*(n,k,\varepsilon)
\le C\min\{n,kL^2/\varepsilon+kL/\varepsilon^2\}.\]

The upper bound is nonadaptive and returns a proper HODLR approximation
with the original fixed-input success guarantee. Corollary 1.2 gives
\(q_*=\Theta(n)\) when \(\varepsilon\le\sqrt{k/n}\). A separate
\href{https://github.com/ajt60gaibb/OpenProblemsInNLA/blob/main/references/holden-re03-2026-09-13/independent-review.md}{independent
Codex AI-agent informal audit} passed these partial results, including
the adaptive two-sided lower bounds.
\href{https://github.com/ajt60gaibb/OpenProblemsInNLA/blob/main/references/holden-re03-2026-09-13/README.md}{Submission,
verified affiliation, sources and reproduction limits}.

\textbf{Still open:} determine the general joint rate up to universal
constants. The uncapped upper expression is \(L\) times the lower
expression, so these results do not solve the full original target. The
problem remains in the open count. The available five-case assembly
smoke suite passed; the historical continuation verification suite is
incomplete in the supplied archive. No Lean verification, external human
peer review or historical novelty claim is asserted.

\newpage

\subsection{Context and notation}\label{context-and-notation}

Use exact real arithmetic, with comparisons, standard Gaussian sampling,
and exact SVDs available as primitives; charge an SVD of an
\(a\times b\) matrix \(O(ab\min(a,b))\) operations. This states the
idealized arithmetic model used here, rather than a finite precision or
bit complexity claim. A matrix--vector query returns either \(Av\) or
\(A^\mathsf Tv\) for one chosen real vector \(v\); both types count
toward the total. Randomized guarantees are for every fixed input, with
probability or expectation over the algorithm's randomness.

\subsection{Problem statement}\label{problem-statement}

Let \(n=2^Lk\) with integers \(k\ge1,L\ge2\), and \(0<\varepsilon<1/2\).
Define \(\mathcal H_{n,k}\) recursively: a matrix of order at most \(k\)
is unrestricted; at a larger node, the two equally sized diagonal blocks
must satisfy the same definition, and each off-diagonal block must have
rank at most \(k\).

Let \(q_*(n,k,\varepsilon)\) be the smallest worst-case number of
queries with which a randomized adaptive algorithm, given only oracle
access to an arbitrary \(A\in\mathbb R^{n\times n}\), returns
\(B\in\mathcal H_{n,k}\) such that

\[\Pr\!\left[
\|A-B\|_F\le(1+\varepsilon)
\min_{H\in\mathcal H_{n,k}}\|A-H\|_F
\right]\ge0.99\]

for every \(A\). Here only oracle calls are charged: arbitrary
measurable computations between queries are allowed. This is an
information complexity question.

\subsubsection{Question}\label{question}

Determine \(q_*(n,k,\varepsilon)\) up to universal multiplicative
constants, including its joint dependence on \(k,L,\varepsilon\).

\subsection{Reference}\label{reference}

Tyler Chen, Feyza Duman Keles, Diana Halikias, Cameron Musco,
Christopher Musco, and David Persson,
\href{https://arxiv.org/pdf/2407.04686v2}{\emph{Near-optimal
hierarchical matrix approximation from matrix-vector products}}, SODA
2025, 2656--2692, Definition 1.1, Problem 1.1, Theorems 1.1--1.2, end of
§6, and §8.

The known upper bound is \(O(\min\{n,kL^4/\varepsilon^3\})\). The source
proves \(\Omega(kL+k/\varepsilon)\) in its regime
\(n\ge c_0k/\varepsilon\) for a sufficiently large absolute \(c_0\);
this restriction matters near the full-recovery threshold.

\subsection{Status evidence}\label{status-evidence}

The latest arXiv version is v2, October 24, 2024. Searches for ``HODLR
query complexity 2026'' and ``HODLR approximation 2026 lower bound''
found no matching improvement. Musco's
\href{https://app.icerm.brown.edu/assets/568/10574/10574_5858_Musco_020420261630_Slides.pdf}{February
2026 ICERM slides}, numbered slide 14, retain the stated upper bound. No
later explicit open-status confirmation was located.

\subsection{Audit --- 2026-09-10}\label{audit-2026-09-10}

Rechecked \href{https://arxiv.org/html/2407.04686v2}{Chen et al.,
Theorems 1.1--1.2 and §8}. The gap in logarithmic and accuracy factors
remains. HODLR query-complexity and later hierarchical-approximation
searches found no matching characterization, including the
near-full-recovery regime.
\end{document}
